% !TeX root = main.tex

\section{INTRODUCTION}

Wall-climbing robots have been developed using mechanical, magnetic, pneumatic, electrostatic, and biomimetic adhesion mechanisms. Although these approaches have enabled locomotion on vertical or inverted surfaces, many systems remain specialized to specific surface materials, environmental conditions, or motion modes. Soft climbing robots address part of this limitation by combining compliant structures with controllable adhesion, allowing crawling, climbing, turning, and, in some cases, transition between horizontal and vertical planes~\cite{Gu2018SoftWallClimbing,Zhang2022Inchworm,Chen2023Flipping,Pang2022Microrobot}. In particular, inchworm-inspired soft robots using pneumatic artificial muscles, negative-pressure suckers, or electroadhesive pads have demonstrated multimodal locomotion and surface-transition capability. However, these systems commonly rely on body contraction or bending patterns optimized for inchworm-like motion, rather than on a lizard-like morphology that combines legged adhesive contact with an actively deformable trunk and tail.

Bioinspired legged climbing robots provide another route toward controllable wall locomotion. Lizard- and gecko-inspired robots have shown that distributed adhesive feet can improve contact stability, maneuverability, and adaptability on vertical surfaces~\cite{Nishad2022Lizard}. These systems demonstrate the importance of legged contact distribution and peeling-based attachment/detachment strategies for climbing. Nevertheless, the body and tail of such robots are often rigid, passive, or treated mainly as supporting structures. As a result, their trunk-tail morphology is rarely used as an active soft mechanism for producing distinct deformation modes required for planar steering and transition between intersecting surfaces.

A related body of work has investigated multi-chamber pneumatic actuators and robotic tails. Three-chamber pneumatic actuators are widely used because they can generate omnidirectional bending with relatively simple pneumatic inputs, although their design is often based on empirical iteration~\cite{Lamping2022ModularActuator}. In parallel, robotic tails have been shown to improve stabilization, maneuverability, terrain adaptation, and climbing robustness in legged robots~\cite{Saab2018RoboticTails,Buckley2023TailStiffness,Yang2023TailControl}. However, most robotic tail studies have focused on rigid or articulated legged robots, while most soft multi-chamber actuators are developed as standalone manipulators or inchworm-like bodies. Therefore, the integration of a modeled multi-chamber soft body-tail actuator into an adhesive lizard-inspired climbing robot remains insufficiently explored.

Motivated by this gap, this paper presents a lizard-inspired pneumatic soft robot that combines adhesive legged locomotion with a coupled three-chamber body-tail actuator. The proposed body-tail actuator consists of two connected soft bending segments that can generate asymmetric deformation because of their pneumatic coupling and inlet configuration. This architecture allows the robot to produce horizontal S-shaped bending for planar steering and vertical U-shaped bending for wall-transition support. Unlike inchworm-inspired soft climbers, the proposed robot uses a lizard-like morphology in which the legs provide distributed adhesive contact while the soft body-tail structure actively contributes to posture control and maneuvering.

The main contributions of this work are as follows:
\begin{enumerate}
    \item The design and fabrication of a lizard-inspired pneumatic soft robot integrating soft adhesive legs with a coupled three-chamber body-tail actuator.
    \item A simplified two-segment constant-curvature model for characterizing asymmetric body-tail bending in horizontal and vertical planes.
    \item Experimental validation of horizontal S-shaped bending, gravity-loaded vertical U-shaped bending, forward locomotion, steering, wall climbing, and wall-transition behaviors.
\end{enumerate}
```